\section{Classical inputs and their formal proofs}\label{sec:classical-inputs}
We record two inputs whose formal implementations are supplied in the
artifact. The first uses the rotation-and-deletion argument attributed to
Woodall~\cite{Woodall1972}; see also~\cite[Theorem 2.7.1]{Gunderson2022}.
The second is a lower-tail estimate proved from finite product weights,
the FKG inequality~\cite{FKG1971}, and a Laplace-transform argument.
These are implementations of classical methods, not claims of new
extremal or correlation inequalities.

\begin{lemma}[Erd\H{o}s--Gallai cycle bound]\label{lem:eg-input}
Let $G$ be a finite simple graph on $n$ vertices and let $k\ge2$ be an
integer. If every simple cycle has length at most $k$, then
\[
                         2\abs{E(G)}\le k(n-1)
\]
when $n\ge1$. For $n=0$ both sides of the corresponding natural-number
statement are zero, with the predecessor interpreted as $n\mathbin{\dot-}1$.
\end{lemma}
\begin{proof}
We use strong induction on $n$. The edgeless case is immediate. Otherwise
choose a longest path $P=v_0\cdots v_\ell$, maximizing $d=\deg(v_0)$ among
the starting degrees of all longest paths. Then $\ell\ge1$. All neighbors
of the first vertex of any longest path lie on that path, since an outside
neighbor would extend it.

Define
\[
 W=\{v_i:0\le i<\ell,\ v_0v_{i+1}\in E(G)\},\qquad
 X=\{v_j:0\le j<\min\{k,\ell+1\}\}.
\]
The predecessor operation is a bijection from $N(v_0)$ to $W$, so
$\abs W=d\ge1$, and $v_\ell\notin W$. For each $v_i\in W$, the rotation
\[
             v_i v_{i-1}\cdots v_0v_{i+1}\cdots v_\ell
\]
is another longest path with the same vertex set. Consequently
$\deg(v_i)\le d$ and $N(v_i)\subseteq V(P)$.

A chord $v_0v_j$ with $j\ge2$ closes a prefix of $P$ to a simple cycle
of length $j+1$, whence $j<k$. The case $j=1$ also has $j<k$.
In particular, $v_i\in W$ implies $i+1<k$, and $W\subseteq X$.
If $v_i\in W$ has a neighbor $v_j$ with $j\ge k$, then $i+1<j$ and
\[
                v_0\cdots v_i v_j v_{j-1}\cdots v_{i+1}v_0
\]
is a simple cycle of length $j+1>k$. Thus every neighbor of $W$ lies in $X$.

Let $a$ count the edges inside $W$ and let $b$ count the edges between
$W$ and $X\setminus W$. The degree bounds and simplicity give
\[
             2a+b\le d^2,\qquad b\le d(\abs X-d).
\]
The $a+b$ edges removed with $W$ therefore satisfy
\[
                     2(a+b)\le d\abs X\le dk.
\]
The remaining graph has $n-d$ vertices, where $1\le n-d<n$, and retains
the same bound on cycle lengths. Applying induction and adding the
deletion cost yields
\[
 2\abs{E(G)}\le dk+k(n-d-1)=k(n-1).
\]
No connectivity hypothesis is used.
\end{proof}

In Lean, the rotations and the degree-maximizing choice are implemented in
\begin{center}\small\nolinkurl{Cycle12/Independent/PathRotation.lean}.\end{center}
The set $W$ and its neighborhood bounds are implemented in
\begin{center}\small\nolinkurl{Cycle12/Independent/WoodallSet.lean}.\end{center}
The counting step is implemented by splitting degree sums across $W$ and
its complement, so internal edges are counted exactly twice in the degree
sum and exactly once in the deletion. The final strong induction is
\mcode{Cycle12.Independent.erdosGallai_cycle}. Passing to a longest cycle
gives the precise long-cycle interface used in the finite completion
argument; its denominator is $n-1$.

\begin{lemma}[Lower tail for a Bernoulli vertex sample]\label{lem:sampling-input}
Let $G$ have $m$ edges, let $D\ge0$, and suppose that
$\deg(u)+\deg(v)\le D+2$ for every edge $uv$.
Independently select each vertex with probability $p\in[0,1]$, and let $X$
count the edges with both endpoints selected. Put
\[
                       \mu=mp^2,\qquad v=mp^2+mDp^3.
\]
If $\mu>0$ and $0\le t\le\mu$, then
\[
                    \Prob(X\le\mu-t)\le
                    \exp\!\left(-\frac{t^2}{2v}\right).
\]
If $\mu=0$, the only required value is $t=0$, for which the bound is the
elementary probability bound by $1$.
\end{lemma}
\begin{proof}
For a subset $S$ of the vertex set, its probability mass is
$w(S)=p^{\abs S}(1-p)^{n-\abs S}$. The identity
\[
                     w(A)w(B)=w(A\cap B)w(A\cup B)
\]
follows by adding the corresponding cardinalities and complement
cardinalities. FKG gives positive association of nonnegative increasing
functions. Shifting an increasing function by its value at the empty set
extends this to real-valued increasing functions. Applying positive
association to $f$ and $M-g$, where $g$ is decreasing and $M\ge\max g$,
gives
\begin{equation}\label{eq:opposite-association}
                 \E(fg)\le (\E f)(\E g)
\end{equation}
for nonnegative increasing $f$ and nonnegative decreasing $g$.
Independence of observables using disjoint coordinates follows by
induction on the finite ground set, conditioning on its last coordinate.

Write $I_e$ for the indicator that edge $e$ survives. For a fixed edge
$e$, split
\[
                         X=I_e+Y_e+Z_e,
\]
where $Y_e$ sums the indicators of the other edges sharing an endpoint
with $e$, and $Z_e$ sums all remaining indicators. There are at most $D$
terms in $Y_e$. Coordinate independence gives, for $s\ge0$,
\[
                  \E(I_e e^{-sZ_e})=p^2\E(e^{-sZ_e}).
\]
For each edge $f$ contributing to $Y_e$, the increasing function $I_eI_f$
has expectation $p^3$; the formal estimate only needs the upper bound
$p^3$. Applying \eqref{eq:opposite-association} and summing gives
\[
                  \E(I_eY_e e^{-sZ_e})\le Dp^3\E(e^{-sZ_e}).
\]
The inequality $e^{-sY_e}\ge1-sY_e$, applied on the event $I_e=1$, now gives
\[
 \E(I_e e^{-sX})\ge e^{-s}(p^2-sDp^3)\E(e^{-sZ_e}).
\]
If $0\le s\le(1+Dp)^{-1}$, the coefficient on the right is nonnegative.
Since $Z_e\le X$, putting $L(s)=\E(e^{-sX})$ yields
\[
 \E(I_e e^{-sX})\ge e^{-s}(p^2-sDp^3)L(s)
                    \ge p^2\bigl(1-s(1+Dp)\bigr)L(s).
\]
For the second inequality, use $e^{-s}\ge1-s$ and then discard the
nonnegative term $s^2Dp^3$. Summation over $e$ gives
\[
                        \E(Xe^{-sX})\ge(\mu-vs)L(s).
\]
Finite summation permits differentiation, with
$L'(s)=-\E(Xe^{-sX})$. Hence
\[
                 s\longmapsto e^{\mu s-vs^2/2}L(s)
\]
is nonincreasing on $[0,\mu/v]$. Its value at $0$ is $1$, so
$L(s)\le e^{-\mu s+vs^2/2}$ there. Finally, the exponential Markov
inequality applied with $s=t/v$ proves the claim. The case $t=0$ is
handled directly by the probability bound by $1$.
\end{proof}

The finite-law calculation is formalized first for a family of distinct
two-element sets, with explicit dependency neighborhoods. The graph
interface instantiates these sets with the endpoints of the original
edges and uses the degree bound in the line graph. The statements
{\small\nolinkurl{Cycle13.vertex_sample_edge_lower_tail}} and
{\small\nolinkurl{Cycle13.finite_graph_janson_restricted}} retain the exact constants
and the finite-sample conventions required later. The proof invokes
Mathlib's FKG theorem and does not assume a separate Janson theorem.
